\documentclass{report}
\usepackage{amsmath,amssymb}
\setlength{\parindent}{0mm}
\setlength{\parskip}{1em}
\begin{document}
\begin{center}
\rule{6in}{1pt} \
{\large Sourish Chakravarty \\
{\bf Characterization of micro-anomalies from macro-scale response}}

252 ENGLEWOOD AVENUE \\ APT2 \\ BUFFALO \\ NY 14214
\\
{\tt sc267@buffalo.edu}\\
Sonjoy Das\end{center}

A stochastic multi-scale based approach is presented in this work to
detect signatures of micro-anomalies \textit{from} macro-level response
variables. In this work, we particularly consider polycrystalline
materials, e.g., Aluminium. Thus, by micro-anomalies, we refer to
micro-cracks of size $10$--$100$~$\mu\mathrm{m}$, while macro-level
response variables imply, e.g., strains, strain energy density of
macro-level structures of typical size often varying at the order of
$10$--$100$~m (e.g., an aircraft wing). For different material and
systems (e.g., sand and geo-mechanical systems), the sizes of micro- and
macro- level scales may, of course, be significantly different. The
micro-anomalies in the context of present work are not discernible by the
naked eye. Nevertheless, they can cause catastrophic failures of
structural systems due to fatigue cyclic loading that results in
\textit{initiation} of fatigue cracks. Analysis of such precursory state
of internal damage evolution, before a macro-crack visibly appears (say,
size of a few cms), is beyond the scope of conventional crack propagation
analysis, e.g., fracture mechanics. The present work is proposed to
address this specific concern, and is an extension of an earlier work by
Das \& Ghanem (2009,2011). In the earlier work, macro-level (continuum)
constitutive properties (e.g., constitutive elasticity tensors) of
heterogeneous materials were constructed within a probabilistic formalism
based on random matrix theory, maximum entropy principle, and principles
of minimum complementary energy and minimum potential energy. The effects
of micro-cracks are now incorporated into the present formulation by
extending the previous work. Distinct differences are observed in the
macro-level response characteristics depending on presence or absence of
micro-cracks. Such stochastic information is used in the framework of
optimization to detect the damaged region, consisting of microcracks,
from experimentally obtained macroscale strain observations. Several
schemes based on both heuristic and traditional optimization techniques
are proposed and tested to detect damage due to microcracks depending on
the availability of information. The proposed work is likely to be useful
in health monitoring of structural systems in, but not limited to,
aerospace, mechanical, and civil engineering applications.


\end{document}
